\subsection{Numerical validation of the tensor algorithm}

The tensor-based formulation was validated against the known analytical
representation transformations of Yamada and
Winnewisser.\cite{Yamada1976}  For quartic centrifugal distortion constants, the
classical transformation between two representations $R$ and $R'$ is
\begin{equation}
\mathbf{D}_{R'} =
Q_{R'}^{-1}(A,B,C)\,Q_R(A,B,C)\,\mathbf{D}_R ,
\end{equation}
where $Q_R$ denotes the Yamada transformation matrix associated with the chosen
representation.

Within the present tensor framework, the same transformation is obtained in a
representation-independent way:
\begin{equation}
\mathbf{D}_{R'}
=
\mathcal P_{R'}\!\left[
P\,\mathcal R_R^{+}(\mathbf D_R)
\right],
\label{eq:quartic_tensor_transform}
\end{equation}
where $\mathcal R_R^{+}$ denotes the pseudoinverse reconstruction of the
pair--pair quartic tensor in the initial representation, $P$ is the permutation
operator associated with the corresponding axis relabeling, and $\mathcal P_{R'}$
denotes the unique forward projection from the transformed tensor to Watson's
quartic constants in the final representation.  In explicit matrix form this
may be written as
\begin{equation}
\mathbf D_{R'}=
\mathbf R_{4,R'}(A',B',C')\,
P\,
\mathbf R_{4,R}^{+}(A,B,C)\,
\mathbf D_R,
\end{equation}
where $\mathbf R_{4,R}$ is the quartic projection matrix for representation $R$.

The equivalence between the analytical and tensorial formulations was verified
using the quartic centrifugal distortion constants of COF$_2$ reported by
Carpenter.  For all representation changes
\[
I_r \rightarrow II_r,\qquad
I_r \rightarrow III_r,\qquad
II_r \rightarrow III_r,
\]
the tensor algorithm reproduces the Yamada transformations to machine
precision, with maximum absolute deviations of order
$10^{-15}$--$10^{-17}$.

For sextic centrifugal distortion constants, for which analogous closed-form
representation transforms are not available, validation was carried out by
round-trip tests.  Starting from sextic constants
\[
\mathbf D^{(6)}=
\bigl(
H_J,H_{JK},H_{KJ},H_K,h_1,h_2,h_3
\bigr)^T,
\]
one performs successive transformations such as
\[
I_r \rightarrow II_r \rightarrow I_r,\qquad
I_r \rightarrow III_r \rightarrow I_r,\qquad
II_r \rightarrow III_r \rightarrow II_r.
\]
For physically admissible sextic tensors, the original constants are recovered
within numerical precision,
\begin{equation}
\max |\mathbf D^{(6)}_{\mathrm{back}}-\mathbf D^{(6)}_{\mathrm{orig}}|
\approx 0.
\end{equation}

These tests confirm the numerical stability of the tensor-based procedure.
For the quartic case it reproduces the classical Yamada transformations exactly,
whereas for the sextic case it provides a self-consistent extension within the
same tensor reconstruction framework.

In practical applications the tensor algorithm also provides useful numerical
diagnostics for comparing different reductions and principal-axis
representations before any new least-squares fit is attempted.  For quartic
constants this role is played by the gauge-sensitive quantity $s_{111}$ and by
the condition number of the corresponding linear transformation.  The former
measures the magnitude of the unitary transformation associated with the chosen
quartic parametrization, whereas the latter quantifies the numerical
amplification of uncertainties in the change of representation.  For sextic
constants the analogous information is provided by the five-dimensional
canonical coordinate vector $\boldsymbol\sigma$, together with the residual
norms and the condition number of the sextic representation transform on the
physical subspace.  Taken together, these diagnostics make it possible to
identify well-conditioned representations and to compare parameter sets coming
from different sources on a common tensorial basis before refining the final
spectroscopic fit.

\section{Results and Discussion}

The tensor reconstruction algorithm was first validated for quartic centrifugal
distortion constants by comparison with the analytical
representation-transform formulas of Yamada and
Winnewisser.\cite{Yamada1976} Using the same rotational constants and quartic
parameters reported in that work, the present tensor procedure reproduces the
analytical results with machine precision for all representation changes.  This
provides a stringent numerical validation of the quartic implementation.

After this validation step, the method was applied to molecules for which
reliable experimental centrifugal distortion constants are available.
Quantum-chemical calculations were then used to assess the performance of
different theoretical levels in predicting the tensor quantities from which the
spectroscopic constants are obtained.

In particular, calculations were performed using two members of the Pisa
Composite Scheme (PCS) family.  The hybrid PCS2 scheme (HPCS2), based on the
B3LYP/6-31+G* level of theory, provides reliable anharmonic contributions to
rotational spectroscopic parameters.  The double-hybrid PCS3 scheme (DPCS3),
based on the revDSD-PBEP86-D3BJ functional combined with the 3F12- basis set,
yields highly accurate harmonic force fields and equilibrium geometries.

In practical applications the two schemes therefore play complementary roles:
HPCS2 provides reliable anharmonic corrections, whereas DPCS3 supplies accurate
harmonic contributions.  Equilibrium rotational constants obtained at the DPCS3
level can be further refined through a simple empirical correction applied to
the bond lengths, leading to the BDPCS3 model, which has been shown to provide
highly accurate equilibrium rotational constants for medium-sized molecules.

Within the present tensor framework, these quantum-chemical calculations do not
simply provide another set of Watson centrifugal distortion constants.  Rather,
they provide access to the underlying tensor quantities from which those
constants are obtained by projection.  This is precisely the point at which the
direct and inverse viewpoints meet.  From the direct side, force fields and
rovibrational couplings lead to reduced quartic and sextic tensors; from the
inverse side, fitted spectroscopic constants can be mapped back to tensor
representatives by pseudoinverse reconstruction.  The two descriptions can
therefore be compared within a common tensorial space before any final choice of
Watson reduction or axis representation is imposed.

In the following sections this framework is applied to representative molecular
systems in order to illustrate both the practical representation transforms and
the tensor-level analysis of centrifugal distortion parameters.  Particular
attention is devoted to the distinction between tensorial content and
representation-dependent parametrization, since the latter may vary
substantially even when the underlying rovibrational information is essentially
the same.  A recurring theme will be that tensor reconstruction allows one to
compare alternative representations on equal footing, assess their numerical
conditioning through quantities such as $s_{111}$ and transformation condition
numbers, and only then decide which representation is most suitable for a
subsequent spectroscopic refit.
